\documentclass[aspectratio=169,11pt]{beamer}

\usetheme{Madrid}
\usecolortheme{default}
\usefonttheme{professionalfonts}
\setbeamertemplate{navigation symbols}{}
\setbeamertemplate{footline}[frame number]

\usepackage{amsmath, amssymb, booktabs}

\title{Fertility, Mortality, Caregiving, And Lifecycle Labor Allocation}
\subtitle{Timing, Persistence, Retirement, And Late-Life Work}
\author{Special Topic 6: Labor Market, Health, and Population -- Week 3}
\date{}

\begin{document}

\begin{frame}
  \titlepage
\end{frame}

\begin{frame}{Central question}
\begin{itemize}
    \item How do fertility timing, mortality risk, aging, and caregiving burdens reshape labor allocation over the life cycle?
    \item The labor objects are hours, participation, earnings, occupation, mobility, savings, retirement timing, bridge jobs, and welfare.
    \item Week 3 discipline: current shocks can move future wages, occupations, savings, retirement, and household specialization.
\end{itemize}
\end{frame}

\begin{frame}{Why this is labor economics}
\begin{itemize}
    \item Fertility affects work because births arrive at specific career stages and change time allocation.
    \item Severe health shocks and mortality risk change household resources, care needs, insurance value, and expected horizons.
    \item Caregiving changes feasible hours, jobs, mobility, retirement timing, and demand for flexibility.
    \item Aging societies make pension rules, older-worker job design, rehiring, and post-retirement work central labor-market margins.
\end{itemize}
\end{frame}

\begin{frame}{Lifecycle labor-allocation framework}
\small
\[
V_t(s_t) =
\max_{h_t,c_t,j_t,a_{t+1},r_t}
u(C_t,L_t,q_t,\ell_t)
+ \beta E_t[V_{t+1}(s_{t+1})]
\]
\begin{itemize}
    \item States: age, health, children, spouse and parent health, assets, pension eligibility, job type, accumulated experience.
    \item Choices: market work, care time, occupation, mobility, savings, retirement, bridge work, post-retirement work.
    \item Shocks: birth, severe illness, death, caregiving onset, pension eligibility, health decline, employer flexibility.
    \item Dynamic point: today's hours, care, and job choices alter tomorrow's wages, options, savings, and welfare.
\end{itemize}
\end{frame}

\begin{frame}{Fertility timing and child penalties}
\small
\begin{tabular}{p{0.28\linewidth} p{0.36\linewidth} p{0.24\linewidth}}
\toprule
Object & Labor channel & Dynamic consequence \\
\midrule
Birth timing & leave, hours, childcare, job continuity & experience path \\
Child penalty & gendered earnings and hours divergence & persistent gaps \\
Occupation choice & flexibility, travel, task demands & wage growth tradeoff \\
Promotion timing & tenure and ladder interruption & slower advancement \\
Specialization & household market-care split & future outside options \\
\bottomrule
\end{tabular}
\vspace{0.3cm}
\begin{itemize}
    \item Separate timing effects from persistent specialization and career-path effects.
\end{itemize}
\end{frame}

\begin{frame}{Family responses to severe health shocks}
\small
\begin{tabular}{p{0.25\linewidth} p{0.36\linewidth} p{0.25\linewidth}}
\toprule
Shock channel & Household response & Identification caution \\
\midrule
Income loss & spouse or adult child works more & anticipation \\
Care demand & family member works less or switches jobs & hidden intensity \\
Mortality risk & savings, insurance, retirement plans change & expected horizon \\
Health capacity & own work declines after illness & severity \\
Liquidity stress & short-run hours and asset responses & insurance access \\
\bottomrule
\end{tabular}
\vspace{0.25cm}
\begin{itemize}
    \item The sign of family labor supply is ambiguous because income replacement and care time push in opposite directions.
\end{itemize}
\end{frame}

\begin{frame}{Caregiving and household reallocation}
\begin{itemize}
    \item Informal care is a labor-market margin: it changes hours, participation, job choice, mobility, and retirement.
    \item Parent care, spouse care, and elder care differ in predictability, intensity, co-residence, distance, and formal-care substitutes.
    \item Paid leave and formal care can preserve attachment by lowering the time cost of family health shocks.
    \item Hidden margins matter: forgone overtime, declined promotion, reduced travel, lower search, and earlier claiming may be missed.
\end{itemize}
\end{frame}

\begin{frame}{Aging, pensions, retirement, and post-retirement work}
\small
\begin{tabular}{p{0.27\linewidth} p{0.36\linewidth} p{0.24\linewidth}}
\toprule
Margin & What moves & Labor interpretation \\
\midrule
Pension eligibility & benefit claiming and exit timing & incentive effect \\
Retirement savings & ability to reduce work or absorb shocks & insurance margin \\
Health decline & capacity and job feasibility & work constraint \\
Flexible work & hours, tasks, schedules, remote work & job-design margin \\
Bridge job / rehire & partial exit, return, part-time work & retirement transition \\
\bottomrule
\end{tabular}
\vspace{0.3cm}
\begin{itemize}
    \item Retirement is often a transition into a different labor arrangement, not simply an absorbing exit.
\end{itemize}
\end{frame}

\begin{frame}{Global aging evidence}
\begin{itemize}
    \item Cross-country evidence helps separate health, preferences, pension rules, employment protection, and care systems.
    \item Population aging changes labor-force attachment, retirement timing, care demand, older-worker job design, and firm adaptation.
    \item The labor question is not demographic drag by itself; it is which work, wage, care, productivity, and retirement margins adjust.
    \item Week 5 will scale up to aggregate demographic change; Week 3 keeps the lifecycle household and retirement margins visible.
\end{itemize}
\end{frame}

\begin{frame}{Research Lab design}
\begin{itemize}
    \item Primary anchor: Fadlon and Nielsen on family labor-supply responses to severe health shocks.
    \item Challenge anchor: Ameriks and coauthors on older-worker flexibility and late-life job design.
    \item Reproduce: estimate a synthetic family health-shock event profile.
    \item Diagnose: classify income replacement, care demand, work-capacity loss, pension incentives, flexibility, and selection.
    \item Transfer: redesign the logic for child penalties, caregiving onset, pension reform, or bridge-job transitions.
\end{itemize}
\end{frame}

\begin{frame}{Research design checklist}
\begin{enumerate}
    \item Name the event: birth, illness, death, care onset, pension eligibility, aging threshold, or retirement transition.
    \item Name the margin: hours, participation, occupation, savings, retirement, rehire, bridge job, or welfare.
    \item State the horizon: immediate disruption, medium-run reallocation, long-run career path, or late-life adjustment.
    \item Separate income replacement, care demand, work capacity, pension incentives, flexibility, and selection.
    \item Say what remains invisible in earnings or employment.
\end{enumerate}
\end{frame}

\begin{frame}{Bridge to Week 4}
\begin{itemize}
    \item Week 3 centers household timing, care burdens, aging, pensions, and retirement transitions.
    \item Week 4 moves inside workplaces: mental health, stress, productivity, absenteeism, presenteeism, retention, and worker welfare.
    \item The common thread is dynamic labor adjustment after health and demographic shocks.
\end{itemize}
\end{frame}

\begin{frame}{Takeaways}
\begin{itemize}
    \item Fertility, mortality, caregiving, and aging are labor-market events because they change work, wages, job choice, savings, and welfare.
    \item Timing and persistence are central: a current shock can alter future career and retirement states.
    \item Caregiving and late-life flexibility are core allocation margins, not side topics.
    \item Global aging matters when it changes institutions and job design around older-worker labor supply and care demand.
\end{itemize}
\end{frame}

\end{document}
